% SPDX-License-Identifier: Apache-2.0
\section{Ziele, Identität und Beziehungen}
\label{sec:core}
\subsection{Das Grundobjekt}
Der normative Kern ist ein endlicher gerichteter Graph
\[
 \mathcal G=(G,C,R),\qquad C,R\subseteq G\times G.
\]
$G$ ist die Menge der Lernziele. $C$ steht für \code{contains},
$R$ für \code{requires}. Es gibt genau diese beiden \emph{ausdrücklich
verfassten} Kantenarten im Zielgraphen. Abgeleitete Relationen, curriculare
Platzierungen, Quellenzuordnungen und Nachweise sind andere Modellschichten.

Ein Paar $(x,y)$ ist gerichtet. Für jede Relation $X$ bedeutet
\[
 X^{-1}=\{(y,x)\mid(x,y)\in X\}.
\]
$(x,y)\in X^+$ gilt genau dann, wenn ein nichtleerer gerichteter $X$-Pfad
von $x$ nach $y$ führt. $X^+$ ist die transitive Hülle;
Pfade der Länge null zählen nicht.
Ein gerichteter Graph ist azyklisch genau dann, wenn
$\forall g\in G:(g,g)\notin X^+$. Die Vereinigung $C\cup R$ ist
\emph{nicht} die Relation, an der pauschal alle DAG-Bedingungen geprüft werden:
Beide Kantenarten haben verschiedene Bedeutungen.

\begin{principle}{Die Pfeilrichtung ist Teil der Definition}
$p\xrightarrow{\texttt{contains}}c$: Der Cluster $p$ enthält das Ziel $c$.
\quad
$g\xrightarrow{\texttt{requires}}p$: Das Ziel $g$ benötigt $p$.
Eine Lernablaufdarstellung zeigt dagegen $R^{-1}$: von der Voraussetzung
zur darauf aufbauenden Kompetenz. Diese Umkehrung muss beschriftet sein.
\end{principle}

\sourcefig[.42\textheight]{01-skill-graph-overview.png}
{Zwei Beziehungen in einem Graphen. Grau gruppiert, Blau zeigt auf
Voraussetzungen. \emph{Velocity} benötigt \emph{Vectors};
\emph{Newton's Second Law} benötigt sowohl \emph{Forces} als auch
\emph{Acceleration}. Die Farben der Knoten sind hier keine Lernzustände.}
{fig:overview}

\subsection{Attribute und stabile Identität}
Die normativen Basisattribute sind totale Abbildungen
\[
 \Id:G\to\mathrm{UUID},\qquad \Title:G\to\Sigma^*,\qquad
 \Weight:G\to\R_{>0}.
\]
$\Id$ ist injektiv: Verschiedene Ziele haben verschiedene IDs.
Die UUID bleibt über Platzierungen, Sprachfassungen und Ansichten hinweg
die kanonische Identität. Kurze Buchstaben in diesem Dokument sind nur
mathematische Bezeichner, keine Aussage zum serialisierten ID-Format.

Ein optionaler $\ShortKey:G\rightharpoonup\Sigma^*$ ist innerhalb einer
logischen Landschaft eindeutig. Mehrere Sprachdateien desselben Ziels
dürfen dasselbe Paar aus ID und ShortKey wiederholen. Die Teilabbildung
$\PhaseCompat:G\rightharpoonup P_{\mathrm{compat}}$ beschreibt dagegen
Kompatibilitäts- oder Anzeigemetadaten, häufig im Feld \code{phase}.
Sie definiert weder kanonische Identität noch die fachliche Lernreihenfolge.
Weitere gebräuchliche Felder ordnet Tabelle~\ref{tab:attributes} ein.

\begin{table}[H]\centering
\begin{tabularx}{\linewidth}{@{}P{29mm}X@{}}
\toprule
Feld & Fachliche Bedeutung außerhalb des minimalen Kerns \\
\midrule
\code{description} & Verständlicher, überprüfbarer Zieltext; typischerweise
\glqq Die lernende Person kann \ldots\grqq. Orientierung hat eine eigene Semantik.\\
\code{area}, \code{tags} & Fachliche Einordnung und Metadaten; ein Tag ersetzt
keine autoritative Entscheidung über Zieltyp oder Geltungsbereich.\\
\code{level} & Schwierigkeit beziehungsweise Anforderungsbereich;
kein curriculares Erwerbslevel und kein Personalisierungslevel.\\
\code{core} & Kernziel versus Erweiterung im jeweiligen Vertrag;
eine Erweiterung kann im gewählten Profil trotzdem verbindliches Ziel sein.\\
\code{weight} & Anteil bei Fortschrittsaggregation; keine Zeitdauer,
keine fachliche Nachweisstärke und keine Bestehenswahrscheinlichkeit.\\
\code{examples} & Referenzen auf Beispiele oder Aufgaben;
Vorhandensein beweist keine Aufgabenqualität oder Kompetenzbeherrschung.\\
\code{resourceLinks} & Bestehendes Referenzfeld unter anderem für eigene
Visualisierungen, offizielle Quellen und Altlinks. Neue optionale externe
Lernmaterialien gehören nach aktuellem Vertrag in getrennte Inhaltspakete.\\
\bottomrule
\end{tabularx}
\caption{Ergänzende Zielattribute und ihre begrenzte Bedeutung.}
\label{tab:attributes}
\end{table}

\subsection{Atomare Ziele und Cluster}
Die \emph{strukturelle} Klassifikation folgt allein aus $C$:
\[
 A=\{g\in G\mid\nexists c\in G:(g,c)\in C\},\qquad K=G\setminus A.
\]
Ein Blatt ist atomar, jeder andere Knoten ein Cluster. Ein Ziel kann mehrere
Eltern haben: Der kanonische \code{contains}-Graph darf eine Polyhierarchie
sein. Explizite Metadaten zur Klassifikation müssen dazu passen.

Strukturelle Atomarität bedeutet nicht automatisch \emph{semantische}
Atomarität. Ein Blatt mit fünf unabhängigen Teilkompetenzen bleibt formal
ein Blatt, kann jedoch die Qualitätsprüfung verfehlen. Orientierung,
Gedächtnisziele und Prüfungen sind besondere pädagogische Arten; sie sind
keine zusätzlichen Strukturtypen und keine Programmabschnitte.

\section{Containment, Voraussetzungen und Vererbung}
\subsection{Hierarchie und atomare Basis}
Für den endlichen azyklischen Containment-Graphen gelten
\begin{align}
 \Anc(g)&=\{a\in G\mid(a,g)\in C^+\},\\
 \Desc(g)&=\{d\in G\mid(g,d)\in C^+\},\\
 \Atoms(g)&=\begin{cases}\{g\},&g\in A,\\\Desc(g)\cap A,&g\in K.\end{cases}
\end{align}
Die Kurzformen $\Anc$ und $\Desc$ entsprechen den Bezeichnern
\code{Ancestors} und \code{Descendants} der Quellspezifikation und ihrer Bilder.
Jeder Cluster besitzt mindestens ein atomares Blatt: In einem endlichen
azyklischen Graphen endet jeder hinreichend lange Pfad. Deshalb ist
$\Atoms(g)$ für jedes Ziel nichtleer. Ein Knoten ohne Kinder ist atomar,
kein leerer Cluster. Bei mehreren Pfaden zählt ein Atom in der Menge nur einmal.

\sourcefig{02-contains-relation.png}
{Direkte und indirekte Zugehörigkeit. Der türkis hervorgehobene Pfad besteht
aus zwei vorhandenen \code{contains}-Kanten; er ist keine dritte Kantenart.
\emph{Mechanics} ist deshalb auch Vorfahr von \emph{Velocity}.}
{fig:contains}
\FloatBarrier

\subsection{Direkte Anforderungen sind konjunktiv}
$(g,p)\in R$ bedeutet, dass $g$ die Voraussetzung $p$ hat. Mehrere
ausgehende Kanten bedeuten \emph{alle} Voraussetzungen, keine Auswahl.
Alternativen brauchen einen eigenen Auswahlvertrag außerhalb dieser Relation.
Die Relation $R$ muss azyklisch sein.

Als kanonisches Ziel der Modellierung wird $R\subseteq A\times A$
bevorzugt. So liegt fachliche Sequenzierung auf den konkreten Kompetenzen.
Eine am Cluster verfasste Voraussetzung ist stärker als eine summarische
Darstellung: Sie gilt durch Vererbung für alle Nachfahren.

\sourcefig{03-direct-requires.png}
{Direkte Voraussetzungen. Auch der gebogene Pfeil ist eine direkte Kante.
Die Kette von Beschleunigung über Geschwindigkeit zu Vektoren ergibt
transitive Erreichbarkeit, keine Containment-Vererbung.}
{fig:requires}
\FloatBarrier

\subsection{Effektive Anforderungen}
Die effektive Relation entsteht aus eigenen und von Vorfahren geerbten
direkten Anforderungen:
\begin{equation}
 (g,p)\in\Reff\iff
 (g,p)\in R\ \lor\ \exists a\in\Anc(g):(a,p)\in R.
 \label{eq:effective}
\end{equation}
Dabei sind
\[
 R_{\mathrm{inherited}}=\Reff\setminus R,\qquad
 \Reff=R\cup R_{\mathrm{inherited}},\qquad
 \Pre_{\mathrm{eff}}(g)=\{p\mid(g,p)\in\Reff\}.
\]
Mehrere Elternpfade addieren Anforderungen als Mengenvereinigung, ohne
Mehrfachzählung. \emph{Vererbung} in \eqref{eq:effective} und
\emph{transitive Hülle} $\Reff^+$ sind verschiedene Operationen.
Die Anforderungen von Voraussetzungen werden durch Pfade geprüft,
nicht durch ein neues Vererbungsprinzip.

Sind alle direkten Anforderungen atomar verfasst, gilt $\Reff=R$:
Containment-Vorfahren sind Cluster und tragen dann keine ausgehenden
Anforderungen. Transitive Ketten bleiben trotzdem relevant.

\sourcefig{04-effective-requires.png}
{Eine Cluster-Anforderung wird vererbt. \emph{Vectors} liegt in diesem
Beispiel außerhalb von \emph{Kinematics}. Würde es auch darin enthalten sein,
erbte es die Anforderung an sich selbst. Durchgezogen heißt direkt verfasst,
gestrichelt heißt geerbt.}
{fig:inheritance}
\FloatBarrier

\section{Gültigkeit, Minimalität und Lernbarkeit}
\label{sec:validity}
\subsection{Verbindliche Bedingungen}
Ein vollständiger Graph ist unter dem normativen Kern genau dann gültig,
wenn die Attribut- und Identitätsbedingungen sowie die folgenden Bedingungen
erfüllt sind:
\begin{align}
 &\forall g\in G:\ (g,g)\notin C^+,\quad (g,g)\notin R^+,
                  \quad (g,g)\notin\Reff^+,\label{eq:acyclic}\\
 &(g,p)\in R\Longrightarrow
      \nexists a\in\Anc(g):(a,p)\in R.\label{eq:localmin}
\end{align}
Zusätzlich gilt \emph{transitive Minimalität}. Für jede direkte Kante
$e=(g,p)\in R$ wird $R'=R\setminus\{e\}$ gebildet und die effektive
Relation $\Reff'$ aus $C,R'$ vollständig neu berechnet. Gefordert ist
\begin{equation}
 (g,p)\notin(\Reff')^+.
 \label{eq:transmin}
\end{equation}
Die direkte Kante muss also für die effektive Erreichbarkeit notwendig sein.
Nur die sichtbare Kante aus $\Reff$ zu löschen wäre falsch: Mit ihr können
auch von ihr abgeleitete Vererbungen entfallen.

Die DAG-Bedingung für $R$ folgt bereits aus derjenigen für $\Reff$, weil
$R\subseteq\Reff$. Sie bleibt als verständliche eigenständige
Autorenprüfung aufgeführt. Unabhängig davon bleibt die Azyklizität von $C$
erforderlich. Identitäten und Referenzen müssen auflösbar sein; unbekannte
Knoten sind keine erfüllten Voraussetzungen.

Auch lokale Minimalität~\eqref{eq:localmin} folgt bereits aus
transitiver Minimalität~\eqref{eq:transmin}: Bei azyklischem $C$ ist jeder
Vorfahr $a$ von $g$ verschieden von $g$. Eine redundante Kante $(g,p)$
zu entfernen lässt daher die Vorfahrenkante $(a,p)$ bestehen;
$(g,p)$ bliebe in $\Reff'$ und somit in $(\Reff')^+$.
Die lokale Bedingung bleibt als einfach prüfbarer Autorenhinweis ausdrücklich genannt.

\sourcefig{05-validity-and-cycles.png}
{Direkte Azyklizität genügt nicht. Aus $A$ enthält $B$ und $A$ benötigt $X$
folgt die geerbte Kante $B\to X$. Kommt $X\to B$ hinzu, entsteht der
effektive Zyklus $B\to X\to B$.}
{fig:cycles}
\FloatBarrier

\subsection{Zusatzprofile und Gegenbeispiele}
Ein Baumprofil kann höchstens einen Elternknoten fordern:
\[
 \forall c\in G:\quad
 \bigl|\{p\in G\mid(p,c)\in C\}\bigr|\leq1.
\]
Dies ist keine Grundbedingung der Polyhierarchie. Ebenso sind Anforderungen
an eigene Vorfahren unerwünscht und in bestimmten Validatorprofilen
verboten. Anforderungen an eigene Nachfahren erzeugen durch Vererbung sogar
eine effektive Selbstabhängigkeit und scheitern bereits am DAG-Vertrag.

\begin{principle}{Ein DAG ist noch kein Beweis der Lernbarkeit}
Seien $G=\{k,a\}$, $C=\{(k,a)\}$ und $R=\{(a,k)\}$.
Die direkten und effektiven Relationen sind azyklisch und minimal.
Aber $\Atoms(k)=\{a\}$: Um $a$ zu beginnen, müsste $a$ schon erfüllt sein.
Die globale Frontier aus leerem Lernzustand ist leer.
Zusätzliche Verbote von Anforderungen an Vorfahren oder ein echter
Abschlusszeuge decken diesen Fehler auf.
\end{principle}

Deshalb müssen drei Aussagen getrennt bleiben: Grundgraph gültig,
gewählte Projektion gültig und Zielmenge aus einem deklarierten Eintrittszustand
vollständig erlernbar. Ein grüner einzelner Validator beweist nicht
automatisch alle drei Aussagen.

\subsection{Didaktische Routen von Orientierung zu Autonomie}
Seien $O\subseteq A$ die Orientierungsanker, $Z\subseteq A$ die
Endziele selbstständiger Leistung und $E_{\mathrm{route}}\subseteq A$
explizit ausgenommene Unterstützungsziele. Die direkte atomare Relation sei
% O and Z avoid collisions with mastery and target sets.
$R_A=R\cap(A\times A)$. Ohne Ausnahmeprofil gilt
$E_{\mathrm{route}}=\varnothing$; Ausnahmen benötigen einen dokumentierten,
maschinenlesbaren Auswahlvertrag. Ein Atom $a\in A\setminus E_{\mathrm{route}}$
ist routenabgedeckt, wenn
\begin{equation}
 \begin{split}
 \exists o\in O,\ z\in Z:\quad
 &(a=o\ \lor\ (o,a)\in(R_A^{-1})^+)\\[-1mm]
 {}\land\;&(a=z\ \lor\ (a,z)\in(R_A^{-1})^+).
 \end{split}
\end{equation}
Die Mengen $O$ und $Z$ heißen in Abschnitt 8.5 der englischen
Quellspezifikation $M$ und $T$. Die abweichenden Buchstaben vermeiden hier
Verwechslungen mit $M_A$ sowie $T_\pi$ und $T_i$.
Die Beschränkung auf $R_A$ verhindert, dass ein Cluster als versteckter
Zwischenknoten einen atomaren Routennachweis liefert. Diese Formel ist
ein Strukturkriterium, kein Beweis didaktischer Angemessenheit.

Konkrete Routenprofile schränken zusätzlich die Schulstufe, erlaubte
semantische Zielarten und die Reihenfolge ein: Orientierung am Einstieg,
fachliches Lernen, bei Bedarf Gedächtnisunterstützung und eine selbstständige
Leistung am Ende. Gedächtnisziele werden dadurch nicht überall Pflicht.
Die Profilregeln \code{CQR-101}, \code{CQR-102}, \code{CQR-103} und
\code{GVR-012}/\code{GVR-013} haben jeweils einen benannten Prüfumfang.
Sie sind nicht als pauschale Zusicherung für jeden Graphen zu lesen.

Eine Übersicht zwischen Clustern kann aus atomaren Routen abgeleitet werden.
Beispielsweise bedeutet eine existenzielle Zusammenfassung $k_2\to k_1$:
\[
 \exists a\in\Atoms(k_2),\ b\in\Atoms(k_1):(a,b)\in R_A^+.
\]
Eine prozentuale Abdeckung braucht darüber hinaus einen erklärten Nenner.
Solche Übersichten erzeugen keine neuen direkten Cluster-Anforderungen.

\section{Lernzustand, Erfüllung und Fortschritt}
\label{sec:state}
\subsection{Erfüllung statt Rückschluss auf Voraussetzungen}
Für die mathematische Progression sei $M_A\subseteq A$ die Menge der
nach dem jeweiligen Zieltyp als erfüllt geltenden atomaren Ziele.
Jedes atomare Lernziel hat genau einen binären Ergebniszustand:
geschafft ($a\in M_A$) oder noch nicht geschafft ($a\notin M_A$).
Teilweise Bearbeitung ist kein dritter Mastery-Zustand.
Die Bezeichnung \glqq mastered set\grqq{} der Spezifikation ist eine
Abstraktion: Bei Orientierung steht Mitgliedschaft für Abschluss, nicht
für fachliche Beherrschung. Bei Gedächtnis- und Prüfungszielen gelten die
jeweiligen besonderen Nachweisregeln.
\begin{equation}
 \Sat(g,M_A)\iff
 \Atoms(g)\ne\varnothing\ \land\ \Atoms(g)\subseteq M_A.
 \label{eq:sat}
\end{equation}
Eine bereits erfolgreiche Anwendung impliziert keine gespeicherte
Beherrschung ihrer Voraussetzungen. $M_A$ wird nicht automatisch
voraussetzungsabgeschlossen. Auch ein als erfüllt angezeigter Cluster
berechtigt nicht zum pauschalen Schreiben von Nachweisen seiner Kinder.

\subsection{Globale Frontier}
Die Menge der nächsten verfügbaren atomaren Ziele lautet
\begin{equation}
 \Frontier(M_A)=\left\{g\in A\setminus M_A\;\middle|\;
 \forall p\in G:\ (g,p)\in\Reff^+\Rightarrow\Sat(p,M_A)\right\}.
 \label{eq:frontier}
\end{equation}
Sie fordert sämtliche transitiv erreichbaren Voraussetzungen. Ein teilweise
bearbeitetes Ziel ist nicht schon deshalb erfüllt. Ein Cluster als
Voraussetzung wird über seine atomare Basis geprüft.

Für Navigation kann getrennt eine Clusterverfügbarkeit definiert werden:
\begin{equation}
 \begin{split}
 \Frontier_K(M_A)=\{k\in K\mid&\ \Atoms(k)\ne\varnothing,
 \ \neg\Sat(k,M_A),\\
 &\forall p\in G:\ (k,p)\in\Reff^+\Rightarrow\Sat(p,M_A)\}.
 \end{split}
\end{equation}
Sie ersetzt die atomare Lernfrontier nicht. Eine Oberfläche kann darüber
hinaus verbindliche Zugangs- und Auswahlregeln anwenden, etwa den
Orientierungseinstieg. Die mathematische Frontier ist deshalb nicht
automatisch identisch mit jedem angezeigten UI-Menü.

\sourcefig{06-available-next-goals.png}
{Mit $M_A=\{\text{Vectors},\text{Forces}\}$ ist nur \emph{Velocity} als
nächstes verfügbar. \emph{Acceleration} wartet auf \emph{Velocity};
\emph{Newton's Second Law} zusätzlich auf \emph{Acceleration}.
Die Cluster sind keine direkt auszuwählenden atomaren Lernschritte.}
{fig:frontier}
\FloatBarrier

\subsection{Aggregation ohne Mehrfachzählung}
Sei $S\subseteq A$ eine fest benannte Menge von Zielatomen,
$w(a)>0$ ihr Fortschrittsgewicht und
$m(a)=\mathbf{1}_{\{a\in M_A\}}\in\{0,1\}$ ihr binärer
Mastery-Indikator. Der gewichtete Fortschritt ist
\begin{equation}
 P(S)=\frac{\sum_{a\in S}w(a)m(a)}{\sum_{a\in S}w(a)}\qquad(S\ne\varnothing).
\end{equation}
Für eine leere Zielmenge gibt es keinen solchen Quotienten; die Anzeige
braucht \glqq keine Ziele im gewählten Umfang\grqq{} statt erfundener 100\,\%.
Für mehrere Cluster wird $S$ als Vereinigung ihrer Zielatome gebildet.
Cluster und Kinder dürfen nicht zusätzlich gemeinsam gezählt werden.
Ein Fortschritt von beispielsweise 50\,\% bezeichnet einen Anteil
geschaffter Ziele im Umfang, keine halbe Beherrschung eines einzelnen Ziels.

Bei der Konvention $w(a)=1$ entspricht das Clustergewicht der Zahl seiner
eindeutigen atomaren Nachfahren. Das ist eine zulässige konkrete
Autorenkonvention, keine Forderung gleicher Gewichte für jede Domäne.
Kern- und Erweiterungsanteile können getrennt ausgewiesen werden; der
verbindliche Personal-Curriculum-Zielumfang entscheidet darüber, was zum
Abschluss gehört. Eine aggregierte Zahl ist weder eine kalibrierte
Wahrscheinlichkeit noch eine Schulnote. Ein Notenmodell bräuchte eine
gesonderte, begründete Abbildung.

\section{Filter, Fokus und globale Diagnose}
\label{sec:filters}
\subsection{Filter als Prädikat und kompilierte Anwendbarkeit}
Ein Filter ist zunächst nur $F:G\to\{0,1\}$ mit
$G_F=\{g\mid F(g)=1\}$. Für dimensionale Anwendbarkeit sei $D$
die Dimensionsmenge und $V_d$ das Vokabular von Dimension $d$. Eine
formal typisierte Menge partieller Zuordnungen ist
\[
 \mathcal V_D=\bigsqcup_{J\subseteq D}\prod_{d\in J}\Pow(V_d),\qquad
 \Applicability:G\rightharpoonup\mathcal V_D.
\]
Die disjunkte Vereinigung hält den jeweiligen Definitionsbereich $J$ fest.
Eine fehlende Dimension ist im allgemeinen Metadatenfilter uneingeschränkt;
eine vorhandene Dimension mit leerer Wertemenge lässt dagegen keinen
konkreten Auswahlwert zu. Nur die Abfrage \code{ALL} ignoriert auch diese
Einschränkung. Die Bedeutung fehlender Metadaten ist keine
Erlaubnis, eine ungeklärte persönliche Schulstufenauswahl als vollständig
aufgelöst zu behandeln.

Für eine aktive Auswahl $Q$ kann die Filterprüfung dimensionsweise erfolgen:
\[
 F_Q(g)=1\iff\bigwedge_{d\in D_{\mathrm{aktiv}}}
 \left\{
 \begin{array}{l}
 Q(d)=\mathrm{ALL},\ \text{oder}\ g\notin\operatorname{dom}(\Applicability),\ \text{oder}\\
 d\notin\operatorname{dom}(\Applicability(g)),\ \text{oder}\\
 Q(d)\in \Applicability(g)(d).
 \end{array}\right.
\]
Die Fälle werden nur ausgewertet, wenn die jeweilige Abbildung definiert
ist. \code{ALL} ist ausschließlich ein Abfragewert und wird nicht als
Geltungsbereich gespeichert. Explizite \code{applicabilityOverrides}
dokumentieren begründete Ausnahmen; die Laufzeit verwendet weiterhin die
kompilierte Anwendbarkeit. Eine Sichtbarkeitsentscheidung beweist noch keine
Quellenabdeckung.

\subsection{Erst global vererben, dann einschränken}
Die induzierten Relationen sind
\[
 C_F=C\cap(G_F\times G_F),\quad R_F=R\cap(G_F\times G_F),\quad
 \Reff|_F=\Reff\cap(G_F\times G_F).
\]
Die effektiven Voraussetzungen werden \emph{vor} der Filterung berechnet.
Sonst könnte eine vom ausgeblendeten Vorfahren geerbte Anforderung
verschwinden, obwohl Ziel und Voraussetzung beide im Fokus liegen.
Die atomare Klassifikation wird ebenfalls nicht aus dem gefilterten Baum
neu erfunden: $A_F=A\cap G_F$. Ein Cluster wird durch das Ausblenden
seiner Kinder nicht zu einem neuen atomaren Ziel.

\subsection{Normalmodus und strikter Modus}
Im Normalmodus (\emph{optimistic}) gilt die relative Basis
\[
 \Atoms_F(g)=\Atoms(g)\cap G_F,\qquad
 \Sat_F(g,M_A)\iff\Atoms_F(g)\subseteq M_A.
\]
Die leere relative Basis blockiert hier nicht. Das bescheinigt keine globale
Kompetenz und verändert keinen gespeicherten Lernzustand. Die Frontier ist
\begin{equation}
 \begin{split}
 \Frontier_{\mathrm{opt}}(M_A,F)=\{g\in A_F\setminus M_A\mid
 \forall p\in G_F:\quad\\
 (g,p)\in(\Reff|_F)^+\Rightarrow\Sat_F(p,M_A)\}.
 \end{split}
\end{equation}
Im strikten Modus bleiben die Kandidaten im Fokus, die Anforderungen
werden jedoch global geprüft:
\begin{equation}
 \Frontier_{\mathrm{strict}}(M_A,F)=
 \{g\in A_F\setminus M_A\mid
 \forall p\in G:\ (g,p)\in\Reff^+\Rightarrow\Sat(p,M_A)\}.
\end{equation}
Somit gilt für dieselben Eingaben
\[
 \Frontier_{\mathrm{strict}}(M_A,F)\subseteq
 \Frontier_{\mathrm{opt}}(M_A,F).
\]
Jeder eingeschränkte Pfad ist ein globaler Pfad; globale Erfüllung impliziert
die relative Erfüllung. Der umgekehrte Schluss gilt nicht.

Auch \emph{Hülle und Einschränkung} sind nicht vertauschbar. Bei
$g\to x\to p$ und $G_F=\{g,p\}$ existiert kein Pfad in
$(\Reff|_F)^+$. Die globale Hülle enthält trotzdem $(g,p)$.
Normalmodus schließt hier die über $x$ laufende Kette aus; strikter Modus
behält sie bei. Das ist die definierte Scope-Semantik.

\sourcefig{07-filter-modes.png}
{Identischer Fokus, unterschiedlicher Prüfbereich. \emph{Vectors} liegt
außerhalb; $M_A=\varnothing$. Normalmodus ignoriert die orange markierten
Anforderungen und erlaubt \emph{Velocity}; strikter Modus lässt keines der
beiden Atome zu. Die Bildverweise 11.3 und 11.4 nennen Abschnitte der
englischen Quellspezifikation.}
{fig:filters}
\FloatBarrier

\subsection{Fehlende Voraussetzungen und fehlende Atome}
Die globale Diagnose trennt benötigte Knoten und noch offene Zielatome:
\begin{align}
 \Missing(g,M_A)&=\{p\in G\mid(g,p)\in\Reff^+\land\neg\Sat(p,M_A)\},\\
 \MissingAtoms(g,M_A)&=
 \left(\bigcup_{p\in \Missing(g,M_A)}\Atoms(p)\right)\setminus M_A.
\end{align}
Für beide Mengen $U$ gelten
$U_{\mathrm{in}}=U\cap G_F$ und $U_{\mathrm{out}}=U\setminus G_F$.
Liegt ein benötigter Cluster im Fokus, sein einziges Atom $b$ aber außerhalb,
kann $\Missing_{\mathrm{out}}=\varnothing$ und dennoch
$\MissingAtoms_{\mathrm{out}}=\{b\}$ gelten. Die Clusterliste allein
verschweigt diese externe Kompetenzlücke.

Die Diagnose folgt der globalen Relation. Sie ist weder die Blockerliste
des Normalmodus noch eine Frontier und erzeugt keine Kanten.

Ein optionaler Prüfbereich $S\subseteq G$ ergibt
\begin{equation}
 \Frontier_{S}(M_A,F)=\{g\in A_F\setminus M_A\mid
 \forall p\in S:\ (g,p)\in\Reff^+\Rightarrow\Sat(p,M_A)\}.
\end{equation}
$S=G$ ergibt den strikten Modus, $S=G_F$ im Allgemeinen \emph{nicht} den
Normalmodus: Pfade und Erfüllung bleiben global. Dieses abgeschwächte
strikte Profil ist ein Modellvorschlag, keine aktivierte dritte Betriebsart.

\noindent\small Quellen dieses Teils:
\srcref{docs/concept/skill-graph/graph-definition.md}{Graph Definition};
\srcref{docs/qa-ci/graph-validation-rules.md}{Graph Validation Rules};
\srcref{AGENTS.md}{Projektkonventionen}.\normalsize
